
\chapter{СБП}\label{ch:sbp}
Система быстрых платежей (СБП)~--- сервис Банка России и~НСПК для
онлайн-переводов и~платежей между счетами в~банках-участниках.
За кнопкой \enquote{карта} стоит схема с~авторизацией, клирингом
и~отдельным расчётом. \highlight{За кнопкой \enquote{СБП} стоит перевод
со~счёта на~счёт, который становится окончательным почти сразу.}

\section{Архитектура: трёхуровневая схема}\label{sec:sbp-architecture}

В~СБП нет карточных ролей: ни~эмитента с~кредитным лимитом,
ни~эквайера с~карточным риском продавца. Вместо двухтактной логики
авторизации и~клиринга СБП устроена по~уровням.

Первый уровень~--- Банк России, через платёжную систему Банка России
(ПС~БР) проходит окончательный расчёт между
банками-участниками~\cite{cbr-ps}. Второй~--- Национальная система
платёжных карт (НСПК; см.~гл.~\ref{ch:mir}), она ведёт операционную
инфраструктуру СБП, вокруг неё построены клиентские сервисы системы
и~правила Операционного платёжного клирингового центра СБП (ОПКЦ
СБП)~\cite{sbp-faq-general,cbr-sbp-actions}. Третий уровень~---
банки-участники, через которые клиент отправляет и~получает перевод.

Если разложить типичный C2B-платёж (consumer-to-business, оплата
покупателем в~пользу продавца; далее в~главе используются типовые
направления денежных потоков~--- C2C между физическими лицами,
B2C от~бизнеса физическому лицу, B2B между юридическими лицами,
Me2Me между своими счетами одного клиента в~разных банках) на~шаги,
последовательность выходит иной, чем у~карты. Продавец создаёт
платёжную сессию (объект на~стороне продавца или банка, связывающий
заказ, попытку оплаты и~итоговый статус) и~показывает динамический
QR-код (quick response code; см.~гл.~\ref{ch:qr-wallets}), кнопку
или ссылку, после чего покупатель уходит в~банковское приложение
или СБПэй и~подтверждает списание с~выбранного счёта. Банк плательщика
передаёт операцию в~операционную инфраструктуру СБП, перевод проходит
расчётную систему Банка России, банк получателя зачисляет деньги,
а~продавец получает итоговый статус
операции~\cite{cbr-ps,cbr-sbp-actions,sbp-faq-business,sbpay-app}.

В~карточной схеме между одобрением и~деньгами есть отдельная стадия
клиринга. В~СБП этот разрыв короче, но экранное подтверждение
не~заменяет исполненный перевод. Пока деньги не~зачислены
и~не~пришёл финальный статус, заказ остаётся в~состоянии ожидания.

\begin{figure}[tbp]
  \centering
  \includefiguresvg[width=\linewidth]{assets/figures/ch18-sbp/sbp-c2b-flow}
  \caption{C2B-платёж по~СБП: путь операции и~точка
  финальности.}\label{fig:sbp-c2b-flow}
\end{figure}

Карточная система работает по~схеме \enquote{сначала авторизация,
потом клиринг}. В~СБП центральное событие~--- подтверждение исполнения
перевода (а~не отдельная авторизация). В~доменной модели нужен другой
набор состояний: инициация, ожидание результата, исполнение, отказ
и~отдельная операция возврата.

\highlight{При оплате картой продавец получает обещание эмитента,
окончательный расчёт приходит позже. При оплате через СБП продавец
ждёт исполненного перевода, и~источник истины о~платеже смещается
от~авторизационного ответа к~факту межбанковского зачисления.}

Если пользовательский интерфейс обновляется быстрее учётной системы
и~заказ помечен как оплаченный раньше, чем перевод подтверждён,
в~сверке появляются расхождения.

\section{Сценарии платежей}\label{sec:sbp-scenarios}

Сценарии СБП различаются тем, кто инициирует перевод, какие данные
нужны на~старте и~где будет проходить сверка. Способ входа~--- номер
телефона, QR-код, ссылка или корпоративные реквизиты~--- для архитектора
вторичен~\cite{sbp-faq-general}.

Перевод от~клиента к~клиенту (C2C)~--- самый короткий сценарий.
Отправитель вводит номер телефона получателя, а~банк отправителя
обращается к~реестру привязок (списку соответствий номеров телефонов
счетам в~банках-получателях, ведётся ОПКЦ СБП) и~определяет, в~каком
банке открыт счёт. Достаточно мобильного приложения банка, номера
телефона получателя и, при необходимости, названия его
банка~\cite{sbp-faq-transfers}. Ограничение здесь одно: номер должен
быть заранее привязан к~счёту в~банке получателя.

Перевод между своими счетами в~разных банках в~СБП называется
Me2Me; он бывает двух видов. \textbf{Me2Me Pull} инициирует
банк-получатель: клиент даёт однократное согласие в~приложении
банка-получателя и~далее одной командой \enquote{стягивает}
средства со~счёта в~банке-отправителе. \textbf{Me2Me Push}
инициирует банк-отправитель: клиент отправляет средства
из~приложения банка-отправителя на~свой счёт в~банке-получателе.
Так пополняют вклад, брокерский счёт, гасят кредит со~счёта
в~другом банке. В~международной карточной практике тот же сценарий
закрывают связкой AFT+OCT (Account Funding Transaction~--- списание
с~карты-источника, Original Credit Transaction~--- зачисление
на~карту-получателя) через Visa Direct или Mastercard Send
(см.~\autoref{sec:visa-direct},~\autoref{sec:mc-send}); в~РФ после марта
2022~года эти продукты для российских участников операционно
недоступны.

Платёж от~клиента к~бизнесу (C2B) проходит через статический
или динамический QR-код. Статический печатают один раз и~используют
многократно; он удобен для вывески, кассы или счёта, но в~сценарии
со~свободной суммой покупатель вводит её сам. Динамический QR-код
создают на~конкретную покупку; он может включать сумму, поэтому
пользователю не~нужно вводить её или перепроверять
вручную~\cite{sbp-faq-business}.

Для продавца разница между статическим и~динамическим QR
сводится к~способу сверки. Статический QR вынуждает сопоставлять
поступивший перевод с~заказом по~сумме, времени, кассе
или дополнительному идентификатору уже после факта оплаты.
Динамический QR лучше подходит к~интернет-магазину и~кассовому ПО,
где платёжная сессия создаётся под конкретную покупку и~связывает
заказ и~перевод~\cite{sbp-faq-business}.

У~маркетплейса один и~тот же C2B-платёж по~СБП обслуживает две
разные модели исполнения заказа. В~FBO (\textit{Fulfillment by
Operator}) товар хранится на~складе площадки, и~она сама готовит
заказ к~доставке. В~FBS (\textit{Fulfillment by Seller}) товар
хранится у~продавца, он сам собирает заказ и~передаёт маркетплейсу
для доставки~\cite{ozon-warehouse-schemes,yandex-marketplace-models}.
На~уровне платёжной системы разницы нет; перевод остаётся
C2B-платежом с~мгновенным зачислением и~отдельной операцией
возврата~\cite{sbp-faq-business}.

Различие в~жизненном цикле заказа: в~FBO хранение и~подготовка
находятся внутри площадки, в~FBS продавец сам отслеживает заказы,
собирает их и~передаёт маркетплейсу для
доставки~\cite{ozon-warehouse-schemes,yandex-marketplace-models}.
Для одной и~той же оплаты по~СБП нужны разные правила состояний
и~сверки.

СБП подходит и~для выплат от~бизнеса физическому лицу (B2C): возвраты,
кешбэк, вознаграждения, выплаты подрядчикам. В~III~квартале 2025~года
в~СБП было совершено 63~млн переводов юридических лиц гражданам
на~сумму 531~млрд рублей~\cite{cbr-sbp-metrics-2025q3}.

Массовые B2C-выплаты требуют загрузки реестра, отдельного статуса
по~каждой строке, повторной отправки ошибочных записей
и~контроля идемпотентности (свойства, при котором повторная отправка
операции с~тем же идентификатором не~приводит к~повторному эффекту).
Продукт управляет не~одной операцией, а~набором связанных переводов.

B2B-сценарий выводит СБП в~корпоративную среду. B2B-перевод требует
формальных реквизитов, назначения платежа и~интеграции с~учётной
системой. Банк России отдельно публикует тарифы для переводов
юридических лиц и~индивидуальных предпринимателей (ИП) в~пользу
юридических лиц и~ИП, а~НСПК описывает B2B-платежи как круглосуточные
онлайн-расчёты с~мгновенным зачислением средств
получателю~\cite{cbr-ps,nspk-sbp-b2b}.

\begin{table}[H]
  \begingroup
  \small
  \setlength{\tabcolsep}{4pt}
  \renewcommand{\arraystretch}{1.2}
  \begin{tabularx}{\textwidth}{
      @{}P{0.11\textwidth}
      P{0.23\textwidth}
      Y
      Y@{}
    }
    \toprule
    \headrow
    \thd{Сценарий} & \thd{Что нужно на старте} &
    \thd{Что подтверждает успех} &
    \thd{Что нужно сверять} \\
    \midrule
    C2C & Номер телефона, при необходимости банк получателя &
    Исполненный перевод на счёт получателя &
    Привязку телефона к~счёту и выбранный банк \\
    Me2Me Pull & Согласие в~приложении банка-получателя на доступ
    к~счёту в~банке-отправителе &
    Исполненный перевод на свой счёт в~банке-получателе &
    Действительность согласия и~доступность счёта-источника \\
    Me2Me Push & Команда из приложения банка-отправителя со ссылкой
    на свой счёт в~банке-получателе &
    Исполненный перевод на свой счёт в~банке-получателе &
    Принадлежность обоих счетов одному клиенту \\
    C2B & QR-код, кнопка, ссылка или привязанный счёт &
    Исполненный перевод в~пользу продавца &
    Связку платёжной сессии и~заказа, а~также возврат как отдельную
    операцию \\
    B2C & Реестр получателей и суммы выплат &
    Реестр исполненных выплат по каждой строке &
    Идемпотентность, частичные ошибки и повторную отправку \\
    B2B & Корпоративные реквизиты и назначение платежа &
    Исполненный перевод и зачисление получателю &
    Связку с~договором, счётом и бухгалтерской системой \\
    \bottomrule
  \end{tabularx}%
  \endgroup
  \caption{Сценарии СБП и~их параметры
  сверки.}\label{tab:sbp-scenarios}
\end{table}

\section{SBPay: мобильный сценарий оплаты}\label{sec:sbp-sbpay}

SBPay (СБПэй)~--- отдельное мобильное приложение для оплаты через СБП.
На~официальной странице сервиса перечислены QR-код, табличка NFC
(near field communication~--- ближняя бесконтактная радиосвязь),
кнопка на~сайте и~платёжная ссылка как поддерживаемые способы
оплаты. Само приложение использует привязанный банковский
счёт~\cite{sbpay-app}.

Пользователь скачивает приложение, выбирает свой банк и~привязывает
счёт через приложение банка или, если это поддерживается, вручную
вводя номер телефона и~номер счёта~\cite{sbpay-app}. После этого
через СБПэй можно оплачивать покупки по~QR-коду, табличке NFC,
кнопке на~сайте или платёжной ссылке.

SBPay~--- отдельный платёжный интерфейс, в~котором банк нужен
для привязки счёта и~части проверок; приложение СБПэй живёт
параллельно банковскому. В~мобильной интеграции заранее решается,
какой путь поддерживается: кнопка СБП в~банковском приложении,
отдельное приложение СБПэй или оба варианта сразу.

В~сценарии с~кнопкой, ссылкой или NFC пользователь уходит во~внешнее
приложение, а~затем возвращается к~продавцу или дожидается фонового
обновления статуса. Поток оплаты проектируется так, чтобы переживать
этот уход: держать платёжную сессию открытой, уметь запрашивать итог
операции повторно и~не~помечать заказ оплаченным по~факту
перехода в~СБПэй или в~банковское приложение.

С~точки зрения интегратора QR-код, табличка NFC, кнопка и~платёжная
ссылка~--- разные способы привести пользователя к~подтверждению
одного и~того же перевода~\cite{sbpay-app,sbp-main}.

\practicenote{%
  Для электронной коммерции проще динамический QR, кнопка СБП
  и~платёжная ссылка~--- они живут внутри сайта или приложения продавца.
  СБПэй нужен там, где вы хотите отдельный мобильный платёжный
  интерфейс поверх СБП.

}
\section{Привязка счёта и~повторные покупки}\label{sec:sbp-subscriptions}

Публично описанный сценарий повторной оплаты в~СБП строится вокруг
привязки счёта. Пользователь один раз выбирает опцию \enquote{Привязать
счёт СБП}, переходит в~приложение банка и~связывает счёт с~продавцом.
После этого можно платить с~привязанного счёта без повторного ввода
реквизитов~\cite{sbp-main}.

Первая онлайн-покупка одновременно проходит как оплата и~как
установление привязки между продавцом и~счётом клиента. Следующая
покупка начинается с~выбора оплаты по~СБП на~сайте или в~мобильном
приложении продавца; повторного ввода платёжных данных уже
не~требуется~\cite{sbp-main}.

Внутри продукта привязку счёта лучше хранить как отдельный объект
согласия между продавцом и~счётом клиента~--- отдельно от~первой
успешной покупки. У~такой привязки собственный жизненный цикл:
запрошена, подтверждена, отозвана, повторная оплата
инициирована. Оплата не~должна неявно следовать из~самой привязки.
Иначе первая же отвязка счёта или отклонение повторной операции
приведут к~смешению состояния согласия и~денежного события.

Снаружи это напоминает card-on-file (сценарий повторных оплат
с~заранее сохранёнными карточными реквизитами), но техническая
модель иная. Продавец не~работает с~PAN (primary account number,
номер карты), карточным токеном и~сроком действия
карты; в~публично описанной модели СБП он опирается на~ранее
привязанный счёт и~дальше строит подтверждаемую пользователем оплату.
Здесь важнее корректно различать момент привязки счёта, момент оплаты
и~обработку результата, чем механически переносить в~СБП карточные
термины MIT (merchant-initiated transaction, операция, инициированная
продавцом) и~CIT (cardholder-initiated transaction, операция,
инициированная держателем карты) по~аналогии.

Публичные материалы СБП подтверждают саму схему привязки и~оплату
с~уже привязанного счёта, но не~раскрывают весь низкоуровневый
протокол между продавцом, банком и~НСПК~\cite{sbp-main}. Точнее
говорить о~привязанном счёте и~отдельной операции, которую всё равно
нужно подтвердить и~учесть по~результату.

По~своей модели СБП ближе к~открытому банкингу (open banking~---
формату, при котором сторонние сервисы по~согласию клиента
инициируют переводы напрямую с~его банковского счёта) и~переводам
со~счёта на~счёт (account-to-account, A2A), чем к~картам. В~обоих
случаях в~центре стоит право инициировать перевод со~счёта клиента,
без карточного реквизита.

\section{Универсальный QR}\label{sec:sbp-universal-qr}

Универсальный QR~--- инициатива НСПК, объединяющая в~одном QR-коде
сразу несколько способов оплаты (по~СБП, банковским сервисом
мгновенной оплаты конкретного банка, цифровым рублём
в~перспективе)~\cite{nspk-universal-qr}. Для покупателя он выглядит
как обычный QR-код. Один и~тот же код в~банковском приложении
автоматически предлагает доступные методы оплаты, а~пользователь
выбирает удобный.

Для продавца Универсальный QR упрощает приём оплат: вместо нескольких
разных кодов и~кнопок на~кассе достаточно одного. Связка \enquote{один
QR~--- много методов} меняет логику сверки: продавец заранее не~знает,
каким способом оплачена конкретная покупка. Учётная система получает
итоговый идентификатор метода в~постфактумном уведомлении и~сверяет
движение денег с~поправкой на~источник перевода.

\section{Подписочные платежи СБП}\label{sec:sbp-subscription}

Помимо привязки счёта (раздел~\ref{sec:sbp-subscriptions}),
НСПК выделяет отдельный сервис \enquote{Подписка СБП}. Однократно
полученное согласие клиента позволяет продавцу инициировать
последующие списания в~пределах согласованных параметров (сумма,
периодичность, срок действия согласия), без участия пользователя
в~каждой повторной операции~\cite{nspk-sbp-subscription}. Внешне
это ближе к~привычной модели card-on-file для подписочного бизнеса,
но техническая основа другая. Вместо токена карты у~продавца
хранится идентификатор согласия СБП с~привязанным счётом клиента.

Для подписочного продукта эта схема прозрачнее привязки счёта.
Согласие явно отделено от~первой оплаты, имеет собственный жизненный
цикл (выдано, активно, отозвано, истекло) и~не~требует от~команды
самостоятельно конструировать объект \enquote{связь со~счётом}
из~истории операций.

\section{Протокол ОПКЦ СБП и~OpenAPI}\label{sec:sbp-protocol}

Технически взаимодействие банка-участника и~продавца с~СБП
строится через OpenAPI (публичный REST-интерфейс, описанный в~формате
OpenAPI Specification) ОПКЦ СБП на~стороне
НСПК~\cite{nspk-sbp-openapi}.\footnote{Полная техническая
спецификация ОПКЦ СБП доступна банкам-участникам и~агентам ТСП
(юридическим лицам, действующим от~имени торгово-сервисного
предприятия в~части приёма платежей по~СБП и~обмена с~ОПКЦ) через
портал \texttt{support.nspk.ru} после подключения; публично опубликованы
Правила информационно-технологического взаимодействия Агента ТСП,
Банка и~АО~\enquote{НСПК} при использовании QR-кодов в~СБП~---
версия~4.0, действует с~12.09.2024.} Ключевые свойства протокола:
асинхронный обмен
сообщениями (создание платёжной сессии и~получение итогового статуса
разделены во~времени), обязательная идемпотентность операций
по~идентификатору запроса, формализованный набор статусов
(\enquote{создан}, \enquote{ожидает подтверждения}, \enquote{исполнен},
\enquote{отказ}, \enquote{отменён}, \enquote{возвращён}) и~явные
поля для связки покупки и~перевода.

Структуру C2B-сессии продавец проектирует вокруг идентификаторов
ОПКЦ: \texttt{qrId} как идентификатор QR-кода или платёжной ссылки,
\texttt{operationId} как идентификатор конкретной попытки оплаты,
отдельный идентификатор возврата. Возврат привязывается именно
к~\texttt{operationId} исходного платежа, а~не к~идентификатору
заказа продавца: один заказ может породить несколько попыток оплаты,
и~для корректного частичного или полного возврата нужно знать,
какой именно перевод сторнируется. Идентификатор кассы и~идентификатор
заказа продавца передаются как пользовательские поля и~не~заменяют
идентификаторы ОПКЦ для целей сверки.

Разработчику стоит сначала освоить часть OpenAPI, открытую
банком-партнёром (она и~определяет реальный набор доступных операций),
затем по~документации НСПК разобрать семантику статусов и~правила
идемпотентности и~только после этого проектировать учётные сценарии
в~продукте.

\section{Финальность, возвраты и~антифрод}\label{sec:sbp-finality}

\highlight{После исполнения перевод нельзя в~одностороннем порядке
отменить ни~банку-отправителю, ни~НСПК. Это свойство самого перевода
(его финальности), а~не интерфейса над ним}~\cite{fz-161}. Если
отправитель ошибся в~сумме или в~номере телефона, вернуть деньги
он сможет с~согласия получателя и~его банка~\cite{sbp-faq-transfers}.
В~C2B-сценарии возврат предусмотрен как новая операция: продавец
инициирует полный или частичный возврат через свой банк, и~деньги
поступают покупателю моментально~\cite{sbp-faq-business}. Встроенного
карточного возвратного платежа в~СБП нет; устройство споров
и~отличия от~карточного цикла подробно разобраны
в~\autoref{sec:disputes-sbp}.

Для учётной системы это означает более строгую машину состояний:
платёжная попытка переходит в~подтверждение исполнения, оттуда~---
в~состояние \enquote{оплачено}; возврат~--- отдельная операция
со~своими статусами. Одного двоичного признака
\enquote{оплачено/возвращено} недостаточно.

\begin{figure}[tbp]
  \centering
  \includefiguresvg[width=\linewidth]{assets/figures/ch18-sbp/sbp-payment-refund-states}
  \caption{Состояния платежа и~возврата в~СБП.}\label{fig:sbp-payment-refund-states}
\end{figure}

Антифрод здесь работает иначе, чем в~карточной схеме: он должен
сработать до~исполнения перевода или в~короткий промежуток сразу
после. По~разъяснению Банка России, при наличии сведений о~получателе
в~базе данных о~переводах без добровольного согласия клиента банк
обязан на~два дня приостановить перевод или отказать в~повторной
операции, в~том числе по~СБП. Если после периода охлаждения
(двухдневной паузы, в~течение которой клиент должен подтвердить
намерение совершить перевод) клиент подтверждает операцию, банк
обязан её исполнить; при соблюдении этой процедуры он не~возвращает
деньги за~перевод на~счёт
злоумышленника~\cite{fz-369-2023,cbr-sbp-faq-cooling-period}.

\section{Сравнение с~карточными схемами}\label{sec:sbp-vs-cards}

Инженерные различия СБП и~карт сведены
в~\autoref{tab:sbp-card-differences}.

\begin{table}[H]
  \begingroup
  \small
  \setlength{\tabcolsep}{4pt}
  \renewcommand{\arraystretch}{1.2}
  \begin{tabularx}{\textwidth}{
      @{}P{0.28\textwidth}
      Y
      Y@{}
    }
    \toprule
    \headrow
    \thd{Параметр} & \thd{Карты} &
    \thd{СБП} \\
    \midrule
    Реквизит на старте & PAN (номер карты), карта или сетевой токен &
    Номер телефона, QR-код, ссылка, кнопка или привязанный счёт \\
    Источник истины о~платеже & Авторизационный ответ, затем
    клиринг и~расчёт &
    Исполненный перевод и подтверждение результата \\
    Финальность денег & После отдельного расчёта &
    Почти сразу после исполнения перевода \\
    Возврат и~спор & Возврат, возвратный платёж, сетевые процедуры &
    Возврат как новая операция, встроенного возвратного платежа нет \\
    Резервирование средств & Возможны холд (резервирование суммы на
    счёте плательщика до завершения клиринга) и предавторизация &
    Естественного аналога карточному холду нет \\
    Повторные покупки & \textit{Stored credential} (сохранённые
    карточные реквизиты) и~правила MIT/CIT &
    Привязанный счёт и отдельное подтверждение результата \\
    География & Международная &
    Внутрироссийская \\
    Предельная сумма & Зависит от карты, банка и~сценария &
    Меньше 1~млн рублей за перевод или платёж \\
    Комиссия для бизнеса & Определяется договором эквайринга & Не
    более 0,4\,\% или 0,7\,\% \\
    \bottomrule
  \end{tabularx}%
  \endgroup
  \caption{Инженерные различия СБП и~карт.}\label{tab:sbp-card-differences}
\end{table}

Карты остаются незаменимыми там, где нужна международная приёмка
и~встроенная карточная процедура спора. СБП выгоднее внутри России,
когда важны мгновенное зачисление и~заранее понятный верхний предел
комиссии для бизнеса~\cite{sbp-faq-transfers,cbr-sbp-c2b-tariffs}.

Продавцу, у~которого большинство чеков укладывается в~лимит одного
перевода, переход на~СБП даёт более простую оплату для клиента
и~меньшую нагрузку на~внутреннюю сверку. За эту простоту приходится
платить другой архитектурой возвратов и~необходимостью самостоятельно
выстраивать процесс разрешения конфликтов.

Карта подходит там, где деньги нужно сначала зарезервировать,
а~уже потом перевести: гостиница, прокат автомобиля, депозит,
предавторизация на~АЗС. СБП сильнее там, где деньги должны быстро
оказаться у~получателя и~заказ закрывается после подтверждённого
перевода.

\section{Тарифы, лимиты и~подключение банков}\label{sec:sbp-tariffs}

Тарифы СБП ниже карточных и~заранее известны. Для переводов между
разными банками комиссия для физического лица отсутствует
до~100~тыс. рублей в~месяц, для переводов между своими счетами
через Me2Me Pull или Me2Me Push (см.~\autoref{sec:sbp-scenarios}) лимит
составляет до~30~млн рублей в~месяц~\cite{sbp-faq-transfers}. Сумма
одного перевода или платежа в~СБП должна быть меньше 1~млн рублей,
банк вправе устанавливать и~более низкие внутренние
лимиты~\cite{sbp-faq-transfers}.

Эти лимиты относятся к~переводам между счетами. Оплата по~QR-коду
или по~кнопке СБП не~расходует лимит на~бесплатные переводы:
в~FAQ системы оплата и~переводы описаны как разные
операции~\cite{sbp-faq-transfers}. Для продуктовой команды слово
\enquote{лимит} здесь скрывает разные ограничения: системный
потолок на~одну операцию, тарифное правило для переводов
и~внутренние ограничения самого банка по~риску и~пользовательскому
пути.

Для приёма платежей бизнеса от~физических лиц (C2B) тариф зависит
от~категории операции. Для социально значимых категорий максимальная
плата для бизнеса ограничена 0,4\,\% от~суммы перевода, для прочих
операций~--- 0,7\,\%~\cite{cbr-sbp-c2b-tariffs}. Тарифные ставки
и~категории C2B-операций устанавливаются решениями Совета директоров
Банка России. Положение Банка России от~24.09.2020~№~732-П
\enquote{О~платёжной системе Банка России} задаёт устройство
самой платёжной системы; ставки и~состав социально значимых категорий
СБП пересматриваются отдельными решениями~\cite{cbr-732p,cbr-sbp-rules}.

Покупатель оплачивает заказ
на~3\,000~руб. в~интернет-магазине электроники. При оплате
по~СБП (несоциально значимая категория) комиссия для бизнеса~---
не~более 0,7\,\%: $3\,000 \times 0{,}007 = 21$~руб. При оплате
картой через эквайринг типичная ставка для интернет-торговли~---
1,5--2,5\,\%; возьмём 2,0\,\%: $3\,000 \times 0{,}02 = 60$~руб.
Разница 39~руб. с~одной транзакции; при 50\,000~операций в~месяц
переход на~СБП экономит порядка 1\,950\,000~руб.
в~год~\cite{cbr-sbp-c2b-tariffs}. Для социально значимых категорий
(аптеки, продукты) потолок 0,4\,\% снижает комиссию до~12~руб.,
увеличивая экономию на~транзакцию до~48~руб.

Число банков-участников СБП ежемесячно публикуется
на~сайте системы~\cite{sbp-faq-general}; к~началу 2026~года
оно превышало 220~кредитных организаций. Для банка-участника
подключение к~СБП формализовано: дополнительное соглашение к~договору
корреспондентского счёта, присоединение к~правилам ОПКЦ СБП и~тестовые
испытания в~инфраструктуре НСПК~\cite{cbr-sbp-actions}. Для бизнеса
маршрут проще~--- нужно обратиться в~банк-участник, поддерживающий
оплату по~QR-коду, и~заключить договор на~приём платежей через СБП;
отдельных соглашений с~НСПК и~Банком России бизнес
не~подписывает~\cite{sbp-faq-business}.

У~интегратора есть выбор. Можно получить от~банка или
провайдера готовый поток с~QR-кодом, кнопкой, ссылкой или оплатой
с~привязанного счёта. Можно построить собственную координацию
процесса: создать динамический QR, связать заказ с~платёжной сессией,
принять обратное уведомление, проверить поступление и~отдельно
запустить возврат, если он нужен. Первый путь быстрее, второй
даёт больше контроля над учётной логикой и~сверкой.

\section{Цифровой рубль и~СБП: разделение функций}\label{sec:sbp-cbdc}

Цифровой рубль похож на~СБП по~интерфейсу: QR-код, банковское
приложение, почти мгновенный результат. Но это иное звено денежной
системы. СБП переводит безналичные деньги между счетами в~коммерческих
банках; цифровой рубль~--- отдельная форма денег, учитываемая
на~платформе Банка России (см.~\autoref{sec:cbdc-architecture}). При
оплате через СБП продавец ждёт перевод между банковскими счетами,
при оплате цифровым рублём~--- движение внутри платформы Банка России.
Одинаковый QR на~экране не~делает эти операции одинаковыми для
реестра обязательств, сверки и~правового
анализа~\cite{cbr-digital-ruble}.

Роль самого Банка России в~этих двух системах разная. В~СБП
он~--- оператор инфраструктуры: ведёт коммутатор, регламент и~расчёт,
но деньги остаются обязательствами коммерческих банков. В~цифровом
рубле он~--- эмитент денежной формы и~оператор реестра кошельков:
обязательство принадлежит Банку России непосредственно, банки лишь
обслуживают клиентский интерфейс. Параллель \enquote{обе системы~---
от~Банка России} верна на~уровне операционного контроля и~неверна
на~уровне того, чьи деньги движутся.
